\documentclass[10pt,a4paper]{article}

\usepackage[utf8]{inputenc}
\usepackage[T1]{fontenc}
\usepackage{lmodern}
\usepackage[top=1.4cm,bottom=1.4cm,left=1.9cm,right=1.9cm]{geometry}
\usepackage{enumitem}
\usepackage[hidelinks]{hyperref}
\usepackage{titlesec}
\usepackage{xcolor}

\definecolor{accent}{HTML}{1F3A5F}

\pagestyle{empty}
\setlength{\parindent}{0pt}
\setlist[itemize]{leftmargin=1.2em, itemsep=2pt, topsep=2pt, parsep=0pt}

\titleformat{\section}{\large\bfseries\color{accent}}{}{0em}{}[\vspace{-0.6em}\color{accent}\rule{\linewidth}{0.5pt}]
\titlespacing*{\section}{0pt}{12pt}{5pt}

% \entry{title}{date}{institution}{location}
\newcommand{\entry}[4]{%
  \textbf{#1} \hfill #2\\
  \textit{#3} \hfill #4\par\vspace{1pt}}

% \talk{date}{event, location}
\newcommand{\talk}[2]{\item #2 \hfill #1}

\begin{document}

% ---------- Header ----------
\begin{center}
  {\LARGE\bfseries Mounir Messaoudi}\\[4pt]
  PhD candidate in Machine Learning, Inria Rennes \quad|\quad \textbf{Available from September 2027}\\[4pt]
  \href{mailto:mounir.mdi@proton.me}{mounir.mdi@proton.me} \quad|\quad
  \href{https://scholar.google.com/citations?user=E_fav18AAAAJ&hl=fr}{Google Scholar} \quad|\quad
  \href{https://github.com/MounirMessaoudi}{GitHub} \quad|\quad
  \href{https://mounirmessaoudi.github.io/}{Website}
\end{center}

% ---------- Research interests ----------
\section{Research Interests}

Robust and usable deep learning methods for 3D biomedical image analysis: detection and segmentation in large, noisy volumes, robustness to acquisition variability and domain shift, and rigorous benchmarking. Applications to cryo-electron tomography, fluorescence microscopy and computational pathology, with tools designed to be used directly by biologists.

% ---------- Education ----------
\section{Education}

\entry{PhD in Deep Learning and Biomedical Image Analysis}{2024 -- 2027 (expected)}
  {Inria, University of Rennes -- SAIRPICO team}{Rennes, France}
Improving deep learning methods for the identification and localization of macromolecules in 3D cryo-electron tomograms; development of a universal particle picker for biologists. Supervisor: Charles Kervrann.

\vspace{6pt}
\entry{M.Sc.\ in Artificial Intelligence, Systems, Data}{2022 -- 2023}
  {Paris Dauphine University -- PSL, with \'Ecole Normale Sup\'erieure and MINES Paris}{Paris, France}

\vspace{6pt}
\entry{M.Sc.\ in Computer Science and Engineering (first year, Erasmus)}{2021}
  {Politecnico di Milano}{Milan, Italy}

\vspace{6pt}
\entry{Engineering Degree, Applied Mathematics and Computer Science}{2017 -- 2022}
  {CY Tech}{Cergy, France}

% ---------- Professional experience ----------
\section{Research and Professional Experience}

\entry{Research Intern in Computer Vision}{03/2023 -- 09/2023}
  {Institut Curie and Centre for Computational Biology (CBIO), MINES Paris}{Paris, France}
\begin{itemize}[label={}, leftmargin=0pt]
  \item Deep learning method to quantify breast cancer aggressiveness from IHC and HES images.
\end{itemize}

\vspace{6pt}
\entry{Data Scientist (work-study, final year of engineering degree)}{09/2021 -- 09/2022}
  {IBM}{Paris, France}
\begin{itemize}[label={}, leftmargin=0pt]
  \item Diagnostics and automatic enhancement of IBM Watson chatbots.
\end{itemize}

% ---------- Publications ----------
\section{Publications}

\begin{enumerate}[label={[\arabic*]}, leftmargin=2.2em, itemsep=6pt, topsep=3pt, parsep=0pt]

\item {\bfseries Designing Affine-Invariant Neural Networks for Photometric Corruption Robustness and Generalization}\\[1pt]
\textbf{M.~Messaoudi}, Q.~Rapilly, S.~Herbreteau, A.~Badoual, C.~Kervrann\\[1pt]
\textit{International Conference on Learning Representations (ICLR)}, 2026 \hfill {\small\textcolor{accent}{\textbf{Poster session}}}\\[1pt]
{\small Links: \href{https://proceedings.iclr.cc/paper_files/paper/2026/file/c9c54ac0dd5e942b99b2b51c297544fd-Paper-Conference.pdf}{[PDF]}}

\item {\bfseries Anisotropy-Aware Strategy for Robust Deep Learning-Based Particle Picking in 3D Cellular Cryo-Electron Tomogram}\\[1pt]
\textbf{M.~Messaoudi}, C.~Kervrann\\[1pt]
\textit{2026 IEEE 23rd International Symposium on Biomedical Imaging (ISBI)}, 2026 \hfill {\small\textcolor{accent}{\textbf{Oral session}}}\\[1pt]
{\small Links: \href{https://ieeexplore.ieee.org/abstract/document/11515672}{[PDF]} \href{https://doi.org/10.1109/ISBI61048.2026.11515672}{[DOI]}}

\item {\bfseries Spatial Point Processes for Cryo-ET Simulation: Validating Particle Placement and Its Effect on Particle Picking}\\[1pt]
\textbf{M.~Messaoudi}, F.~Lavancier, E.~Moebel, F.~Fatmaoui, I.~Tanoz, M.~Eltsov, C.~Kervrann\\[1pt]
\textit{Under review}, 2026

\end{enumerate}

% ---------- Talks ----------
\newpage
\section{Talks and Presentations}

\begin{itemize}[label={}, leftmargin=0pt, itemsep=3pt]
  \talk{11/2026}{\textbf{Invited talk}, Biological Image Analysis Session (BIS 2.0), IGDR, Rennes}
  \talk{07/2026}{CAP \& RFIAP 2026, Montpellier -- Poster session}
  \talk{05/2026}{Life Science PhD Symposium (3rd edition), Rennes -- Poster session}
  \talk{04/2026}{ICLR 2026, Rio de Janeiro -- \textbf{Poster session}}
  \talk{04/2026}{ISBI 2026, London -- \textbf{Oral session}}
  \talk{03/2026}{IABM 2026, Lyon -- Poster session}
  \talk{01/2026}{AI WILD West 2026, Rennes -- Poster session}
\end{itemize}

% ---------- Academic service ----------
\section{Academic Service}

\begin{itemize}[label={}, leftmargin=0pt]
  \item \textbf{Reviewer:} ICLR 2027, AAAI 2027.\end{itemize}

% ---------- Teaching ----------
\section{Teaching}

\begin{itemize}[label={}, leftmargin=0pt]
  \item \textbf{Lectures:} \textit{Introduction to Computer Vision} (20h per year), Master's in Mathematical Engineering (CY Tech) and Master's in Applied Mathematics and Statistics (University of Rennes), 2025 and 2026.
  \item \textbf{Practical sessions:} Python, third-year Bachelor's in Computer Science, ISTIC, University of Rennes (2025).
  %\item \textbf{Tutoring:} Mathematics and sciences tutor for underprivileged high-school and engineering students.
\end{itemize}

% ---------- Skills, languages & interests ----------
\section{Skills and Languages}

\begin{itemize}
\item \textbf{Machine learning}: invariant/equivariant architectures, domain generalization, robustness, localization, detection, classification and segmentation.
\item \textbf{Programming}: Python, PyTorch, Tensorflow, GPU clusters, Git, Linux.
  \item \textbf{Imaging modalities:} 3D cryo-electron tomography, histopathology (IHC, HES), fluorescence microscopy.
  \item \textbf{Languages:} French (fluent), English (fluent).
\end{itemize}

% ---------- References ----------
\section{References}

\begin{itemize}[label={}, leftmargin=0pt]
  \item \textbf{Charles Kervrann} {\small (Inria Rennes)} \quad \href{mailto:charles.kervrann@inria.fr}{charles.kervrann@inria.fr} \hfill \textit{PhD supervisor}
  \item \textbf{Thomas Walter} {\small (Institut Curie / Mines Paris)} \quad \href{mailto:thomas.walter@minesparis.psl.eu}{thomas.walter@minesparis.psl.eu} \hfill \textit{Master's thesis supervisor}
  \item \textbf{S\'ebastien Herbreteau} {\small (ENSAI)} \quad \href{mailto:sebastien.herbreteau@ensai.fr}{sebastien.herbreteau@ensai.fr} \hfill \textit{Co-author}
\end{itemize}

\end{document}
